\paragraph{Floor-entry research.}
When class \(H\) buys and class \(L\) sells after a success and no order recruits, entry is \(\rho\) after every
order, and the purchase and the sale both earn \(\rho\DT m_\lambda-k\) at the conjectured research \(\lambda\). For
\(\rho\DT\ge2k\), the root \(\lambda_Q(\DT)\) of \eqref{eq:a5-floor} exists and is unique, because
the gap \(\rho\DT m_\lambda-k-a\lambda\) falls strictly in \(\lambda\), from
\(\rho\DT/2-k\ge0\) at \(\lambda=0\) to \(-k-a\) at \(\lambda=1\). The root is zero at
\(\rho\DT=2k\) and positive when \(\rho\DT>2k\). The gap rises
strictly in \(\DT\), and its derivative in \(\lambda\) is negative, so \(\lambda_Q\) rises strictly and continuously
in \(\DT\).

For \(\rho\DT(r)>2k\), \emph{floor-entry research} is the tuple with \(\hat\lambda=\min\{\lambda_Q(\DT(r)),\lb\}\),
\(\sigma_H=\delta_1\), \(\sigma_L=\delta_{-1}\), \((\mu(-1),\mu(0),\mu(1))=(m_{\hat\lambda},1/2,M_{\hat\lambda})\),
\(\zeta\equiv0\), prices \(P(q)=t_0+\rho\,w(\mu(q))\), \(\bar\mu(P(q))=\mu(q)\) and truthful bids. Its entry is
\(\rho\).

\emph{Recruiting research at the cap} is the tuple with \(\hat\lambda=\lb\), \(\sigma_H=\delta_1\),
\(\sigma_L(-1)=\hat s_-(r)/\lb\) and \(\sigma_L(0)=1-\hat s_-(r)/\lb\), where \(\hat s_-(r)\) is the low-class sale rate
of Lemma~\ref{lem:a4-low} at the presence \(\lb\) and the prior \(1/2\):
\[
\hat s_-(r)=\min\bigl\{\lb,\ \bar b\,[\rho\DT(r)/k-2]^+\bigr\},\qquad \bar b=(1-\lb)\beta,\qquad
\bar b_0=(1-\lb)(1-2\beta).
\]
Its posteriors are \(\mu(-1)=\bar b/\bigl(2\bar b+\hat s_-(r)\bigr)\),
\(\mu(0)=\bar b_0/\bigl(2\bar b_0+\lb-\hat s_-(r)\bigr)\) and \(\mu(1)=M_{\lb}\). It has \(\zeta(P(1))=1\) and
\(\zeta(P(0))=\zeta(P(-1))=0\), prices \(P(q)=t_0+e(q)\,w(\mu(q))\), \(\bar\mu(P(q))=\mu(q)\) and truthful bids.
Class \(L\) abstains when \(\rho\DT\le2k\) and sells with probability one when \(\rho\DT m_{\lb}\ge k\); in between it
mixes, and its sale just pays the trading cost.

The \emph{floor margin} is
\[
\varphi_\rho(r)=\rho\DT(r)\,[1-\tau(r)]-k-a\lambda_E(r).
\]
Since \(m_{\lambda_E(r)}=1-\tau(r)\), it is the gap \(\rho\DT m_\lambda-k-a\lambda\) of \eqref{eq:a5-floor} at
\(\lambda=\lambda_E(r)\): the return to a full order at the entry floor, at the research that lets a purchase just
reach the threshold, net of the trading cost and of the marginal cost of that research.

\begin{lemma}[Floor-entry research along the path]\label{lem:a5-floor}
Take the example at the prior \(1/2\), with research cost \(a=\reg{res_a}\) and cap \(\lb\approx\reg{res_cap}\), when an
indifferent investor may mix over its best orders and an indifferent high-cost challenger prepares. Let
\(r\in[r_0,r_2]\).
\begin{enumerate}[label=(\alph*),leftmargin=2em]
\item The floor margin \(\varphi_\rho\) rises strictly on \([r_0,r_2]\). Its unique zero \(r_Q\approx\reg{res_rQ}\)
satisfies \(r_N<r_Q<r_U<r_C\). On \((r_N,r_C]\), \(\tau(r)\) exceeds the top posterior \(M_\lambda\) at
\(\lambda=\min\{\lambda_Q(\DT(r)),\lb\}\) exactly on \((r_N,r_Q)\), and equals it exactly at \(r_Q\) and at \(r_C\).
\item Floor-entry research is an equilibrium exactly on \((r_N,r_Q)\) and on \((r_C,r_2]\), with entry \(\rho\). The
root \(\lambda_Q(\DT(r))\) reaches the cap at \(r_U\): \(\lambda_Q(\DT(r))\ge\lb\) exactly when \(r\ge r_U\). On
\([r_N,r_Q]\), \(\lambda_Q(\DT(r))\) stays below the cap and rises strictly from \(0\) at \(r_N\) to
\(\lambda_E(r_Q)\approx\reg{res_lambda_floor_rQ}\) at \(r_Q\). On \((r_C,r_2]\), research is at the cap.
\item Recruiting research at the cap is an equilibrium exactly on \([r_{\lb},r_C]\), with entry
\(\reg{res_entry_cap}\).
\item On \((r_N,r_Q)\), the equilibria are exactly floor-entry research and recruiting research at the cap. On
\([r_Q,r_C]\), recruiting research at the cap is the unique equilibrium. On \((r_C,r_2]\), floor-entry research at
the cap is the unique equilibrium.
\end{enumerate}
\end{lemma}

\begin{proof}
Throughout, \(r\in[r_0,r_2]\), \(\lambda_Q\) stands for \(\lambda_Q(\DT(r))\), and \emph{the gap} is
\(\rho\DT(r)\,m_\lambda-k-a\lambda\), the gap of \eqref{eq:a5-floor} at \(\DT=\DT(r)\), as a function of \(\lambda\).
Recall \(r_N=\fr(2k/\rho)\), \(r_{\lb}=\fr\bigl((2a\lb+k)(1+\eu)\bigr)\),
\(r_U=\fr\bigl((k+a\lb)(1+\eu)/\rho\bigr)\), and the entry ceiling \(r_C\), at which \(\tau(r_C)=M_{\lb}\).

\emph{Step 0: facts along the path.} The spread \(\DT\) and the threshold \(\tau\) rise strictly. Since
\(B_r(1/2)=5-r/4+1/(8r)<c_H\), \(\tau(r)>1/2\), and \(\tau(r)\le\tau(r_2)=\reg{tau_r2}<1\); so \(\lambda_E(r)\in(0,1)\)
is defined and rises strictly. Since \(M_\lambda\) rises strictly in \(\lambda\) and \(M_{\lambda_E(r)}=\tau(r)\),
\(\lambda_E(r)\le\lb\) exactly when \(r\le r_C\), with equality only at \(r_C\). Condition \eqref{eq:a5-U} holds at
\(r_1\), with margin \(\reg{margin_res_bound_r1_B}\), and \(M_{\lb}-\tau(r_1)=\reg{margin_res_tau_r1_B}\) and
\(\tau(r_2)-M_{\lb}=\reg{margin_res_tau_r2_B}\), so \(r_U<r_1<r_C<r_2\). Also \(r_N-r_{\lb}=\reg{margin_res_rN_B}>0\).
Since \(m_{\lb}=1/(1+\eu)\), the gap at \(\lb\), \(\rho\DT(r)\,m_{\lb}-k-a\lb\), is nonnegative exactly when
\(r\ge r_U\), with equality only at \(r_U\). For \(r\ge r_N\), the gap falls strictly in \(\lambda\) and vanishes at
\(\lambda_Q\), so \(\lambda_Q\ge\lb\) exactly when \(r\ge r_U\).

\emph{Step 1: the comparison.} Let \(r\in(r_N,r_C]\). Then \(\lambda_E(r)\le\lb\), and \(m_{\lambda_E(r)}=1-\tau(r)\)
makes the gap at \(\lambda_E(r)\) equal to \(\varphi_\rho(r)\). Since \(M_\lambda\) rises strictly, \(\tau(r)\)
exceeds \(M_\lambda\) at \(\lambda=\min\{\lambda_Q,\lb\}\) exactly when \(\min\{\lambda_Q,\lb\}<\lambda_E(r)\). Because
\(\lambda_E(r)\le\lb\), this is \(\lambda_Q<\lambda_E(r)\), and because the gap falls strictly in \(\lambda\) and
vanishes at \(\lambda_Q\), it is \(\varphi_\rho(r)<0\). In the same way, \(\tau(r)\) equals \(M_\lambda\) at
\(\lambda=\min\{\lambda_Q,\lb\}\) exactly when \(\min\{\lambda_Q,\lb\}=\lambda_E(r)\), which holds exactly when
\(\lambda_Q=\lambda_E(r)\), that is \(\varphi_\rho(r)=0\), or when \(\lb=\lambda_E(r)\le\lambda_Q\), that is \(r=r_C\),
by Step 0 and \(r_C>r_U\).

\emph{Step 2: the floor margin rises.} By the formula for \(d\lambda_E/d\tau\),
\[
\varphi_\rho'(r)=\rho\bigl[\DT'(r)\,(1-\tau(r))-\DT(r)\,\tau'(r)\bigr]
-\frac{a\beta\,\tau'(r)}{\{\beta+(1-2\beta)[1-\tau(r)]\}^2},
\]
with \(\DT'(r)=\tfrac12(1-1/r^2)\) and \(\tau'(r)=3(8r^2-r+2)/[2(20r-r^2-1)^2]\). On \([r_0,r_2]\) the factors
\(\DT'\), \(\DT\), \(1-\tau\) and \(\tau'\) are positive. On a piece \([r',r'']\subset[r_0,r_2]\), \(\DT'\), \(\DT\)
and \(\tau\) rise, \(8r^2-r+2\) rises, and \(\tau'\le T(r',r'')\), since \(20r-r^2-1\) is positive and rises; the
denominator \(\{\beta+(1-2\beta)[1-\tau]\}^2\) of the last term falls as \(\tau\) rises, so on the piece it is at
least its value at \(r''\). Hence, for \(r\in[r',r'']\),
\begin{gather*}
\varphi_\rho'(r)\ge\rho\bigl[\DT'(r')\,(1-\tau(r''))-\DT(r'')\,T(r',r'')\bigr]
-\frac{a\beta\,T(r',r'')}{\{\beta+(1-2\beta)[1-\tau(r'')]\}^2},\\
T(r',r'')=\frac{3(8(r'')^2-r''+2)}{2(20r'-(r')^2-1)^2}.
\end{gather*}
Partition \([r_0,r_2]\) into pieces of width \(\reg{res_floor_step}\). An exact rational computation shows that the
right side of the first line is at least \(\reg{res_floor_dphi_lb}\) on each piece, so \(\varphi_\rho\) rises
strictly on \([r_0,r_2]\).

\emph{Step 3: the zero.} At \(r_N\), \(\rho\DT(r_N)=2k\), so \(\varphi_\rho(r_N)=k[1-2\tau(r_N)]-a\lambda_E(r_N)<0\).
Since \(m_{\lb}=1-M_{\lb}\),
\[
\varphi_\rho(r)=\bigl[\rho\DT(r)\,m_{\lb}-k-a\lb\bigr]+\rho\DT(r)\bigl[M_{\lb}-\tau(r)\bigr]+a\bigl[\lb-\lambda_E(r)\bigr].
\]
On \([r_U,r_C]\), the first bracket is nonnegative, with equality only at \(r_U\), and the other two terms are
nonnegative and vanish only at \(r_C\). Since \(r_U<r_C\), \(\varphi_\rho>0\) on \([r_U,r_C]\). By Step 2,
\(\varphi_\rho\) has exactly one zero \(r_Q\) on \([r_0,r_2]\), with \(r_N<r_Q<r_U\); it is negative on \([r_0,r_Q)\)
and positive on \((r_Q,r_2]\). The floor margin is a rational function of \(r\) with rational coefficients, so its
sign at a rational strength follows from exact rational arithmetic; these signs give \(r_Q\approx\reg{res_rQ}\) and
\(\lambda_E(r_Q)\approx\reg{res_lambda_floor_rQ}\). With Step 1, this proves part (a).

\emph{Step 4: floor-entry research.} Let \(r\in(r_N,r_2]\), so \(\rho\DT(r)>2k\), and take floor-entry research,
with \(\hat\lambda=\min\{\lambda_Q,\lb\}>0\). (E1): \(f_H=(\hat b,\hat b_0,\hat b+\hat\lambda)\) and
\(f_L=(\hat b+\hat\lambda,\hat b_0,\hat b)\) on \((-1,0,1)\) give the stated posteriors, and
\(m_{\hat\lambda}<1/2<M_{\hat\lambda}\). (E2): the stated entry is \(\rho\) after every order, so by the paragraph
on constructed tuples in Appendix~\ref{app:beliefs} the three prices differ, and each price reveals its order.
Suppose \(\tau(r)>M_{\hat\lambda}\). (E3): every posterior is at most \(M_{\hat\lambda}<\tau(r)\), so \(B_r<c_H\) at
every price, and \(\zeta\equiv0\) with no tie. (E4): \eqref{eq:a2-terminal}. (E5): by \eqref{eq:a5-payoffs},
\(U_H(1)=\rho\DT(1-M_{\hat\lambda})-k\) and \(U_L(-1)=\rho\DT m_{\hat\lambda}-k\). Both equal
\(\rho\DT m_{\hat\lambda}-k\), which is at least \(a\hat\lambda>0\), because \(\hat\lambda\le\lambda_Q\) and the gap is
nonnegative on \([0,\lambda_Q]\). With \(U_j(0)=0\) and Lemma~\ref{lem:a5-bounds}(c), the purchase and the sale are
the strict best orders. (E6): truthful bids. (E7): \(\Psi=\rho\DT m_{\hat\lambda}-k\). If \(\lambda_Q<\lb\), then
\(\hat\lambda=\lambda_Q\) and \(\Psi/a=\lambda_Q<\lb\); if \(\lambda_Q\ge\lb\), then \(\hat\lambda=\lb\) and
\(\Psi-a\lb=\rho\DT m_{\lb}-k-a\lb\ge0\). In both cases \(\min\{\lb,\Psi/a\}=\hat\lambda\). Conversely, if
\(\tau(r)\le M_{\hat\lambda}\), then \(B_r(\mu(1))\ge c_H\), and (E3) with an indifferent challenger that prepares
requires \(\zeta(P(1))=1\), so the tuple is not an equilibrium.

Floor-entry research is therefore an equilibrium exactly when \(\tau(r)>M_{\hat\lambda}\) at
\(\hat\lambda=\min\{\lambda_Q,\lb\}\). By part (a), on \((r_N,r_C]\) this holds exactly on \((r_N,r_Q)\). On
\((r_C,r_2]\), \(r>r_U\) gives \(\min\{\lambda_Q,\lb\}=\lb\), and \(\tau(r)>M_{\lb}\), so it holds there, with research
at the cap. On \([r_N,r_Q]\), \(r<r_U\), so \(\lambda_Q<\lb\); it rises strictly and continuously with \(r\), from \(0\)
at \(r_N\), where the gap at \(\lambda=0\) is \(\rho\DT(r_N)/2-k=0\), to \(\lambda_E(r_Q)\) at \(r_Q\), where the gap
at \(\lambda_E(r_Q)\) is \(\varphi_\rho(r_Q)=0\). This proves part (b).

\emph{Step 5: recruiting research at the cap.} Let \(r\in[r_{\lb},r_C]\), and write \(\hat s_-=\hat s_-(r)\). (E1):
\(f_H=(\bar b,\bar b_0,\bar b+\lb)\) and \(f_L=(\bar b+\hat s_-,\bar b_0+\lb-\hat s_-,\bar b)\) give the stated
posteriors, with \(\mu(-1),\mu(0)\le1/2<\tau(r)\le M_{\lb}=\mu(1)\). (E2): the stated entry, \(\rho\) after a sale and
after no order and \(1\) after a purchase, is nondecreasing in the posterior; by the paragraph on constructed tuples,
orders with equal posteriors have equal entry \(\rho\) and share a price, orders with distinct posteriors have
distinct prices, and the belief at each price is the posterior of its orders. (E3): \(B_r(M_{\lb})\ge c_H\), with
equality only at \(r_C\), where the indifferent challenger prepares, and \(B_r<c_H\) at the posteriors at most
\(1/2\). (E4): \eqref{eq:a2-terminal}. (E5): the purchase earns \(\DT m_{\lb}-k\ge\DT(r_{\lb})\,m_{\lb}-k=2a\lb>0\),
so by Lemma~\ref{lem:a5-bounds}(c) it is the strict best order of class \(H\). Class \(L\) earns \(U_L(1)\le-k\),
\(U_L(0)=0\) and \(U_L(-1)=\rho\DT\mu(-1)-k\). If \(\hat s_-=0\), then \(\rho\DT\le2k\) and
\(U_L(-1)=\rho\DT/2-k\le0\). If \(0<\hat s_-<\lb\), then \(2\bar b+\hat s_-=\bar b\rho\DT/k\), and \(U_L(-1)=0\). If
\(\hat s_-=\lb\), then \(\rho\DT m_{\lb}\ge k\), since \(m_{\lb}=\bar b/(2\bar b+\lb)\), and
\(U_L(-1)=\rho\DT m_{\lb}-k\ge0\). In each case \(\sigma_L\) puts weight only on best orders, and an indifferent
investor may mix. (E6): truthful bids. (E7): \(\Psi\ge\tfrac12(\DT m_{\lb}-k)\ge a\lb\), so
\(\min\{\lb,\Psi/a\}=\lb\). A purchase has probability \(\tfrac12(\bar b+\lb)+\tfrac12\bar b=\beta+(1/2-\beta)\lb\), and
only the purchase raises entry above \(\rho\), so entry is \(\rho+(1-\rho)[\beta+(1/2-\beta)\lb]=\reg{res_entry_cap}\).

For \(r>r_C\), \(\mu(1)=M_{\lb}<\tau(r)\), so (E3) requires \(\zeta(P(1))=0\), and the tuple is not an equilibrium. For
\(r<r_{\lb}\), \(r<r_N\) gives \(\hat s_-=0\) and \(U_L(-1)<0\), so \(\Psi=\tfrac12(\DT m_{\lb}-k)^+<a\lb\), because
\(\DT(r)<\DT(r_{\lb})\); (E7) fails. This proves part (c).

\emph{Step 6: every equilibrium above \(r_N\).} Let \(r\in(r_N,r_C]\), and take any equilibrium. If
\(\hat\lambda=0\), every posterior is \(1/2\), entry is \(\rho\) after every order, and \(U_H(1)=\rho\DT/2-k>0\); then
\(\Psi>0\), and Lemma~\ref{lem:a5-joint} gives \(\hat\lambda>0\), a contradiction. So \(\hat\lambda>0\), and
Lemma~\ref{lem:a5-buy} gives \(\sigma_H=\delta_1\), \(\mu(1)=M_{\hat\lambda}\) and \(e(0)=e(-1)=\rho\).
Lemma~\ref{lem:prices}(d) gives \(e(1)=1\) if \(M_{\hat\lambda}\ge\tau(r)\) and \(e(1)=\rho\) otherwise.

If \(M_{\hat\lambda}<\tau(r)\), then \(U_H(1)=\rho\DT m_{\hat\lambda}-k\), which is positive by
Lemma~\ref{lem:a5-buy}. Lemma~\ref{lem:a5-bounds}(a) gives \(\mu(-1)\ge m_{\hat\lambda}\), so
\(U_L(-1)\ge\rho\DT m_{\hat\lambda}-k>0\), and the sale is the unique best order of class \(L\). Then
\(\sigma_L=\delta_{-1}\), \(\mu(-1)=m_{\hat\lambda}\) and \(\Psi=\rho\DT m_{\hat\lambda}-k\).
Lemma~\ref{lem:a5-joint} gives \(\hat\lambda=\min\{\lb,\Psi/a\}\): if \(\hat\lambda<\lb\), then
\(a\hat\lambda=\rho\DT m_{\hat\lambda}-k\), so the gap vanishes at \(\hat\lambda\) and \(\hat\lambda=\lambda_Q<\lb\);
if \(\hat\lambda=\lb\), then the gap at \(\lb\) is nonnegative, and \(\lambda_Q\ge\lb\). So
\(\hat\lambda=\min\{\lambda_Q,\lb\}\). The orders give the posteriors, Lemma~\ref{lem:prices} gives the preparation,
the prices and the beliefs, and the bids are truthful, so the tuple is floor-entry research.

If \(M_{\hat\lambda}\ge\tau(r)\), then \(e(1)=1\) and \(U_H(1)=\DT m_{\hat\lambda}-k\). Class \(L\) can abstain, so
\(\Psi\ge\tfrac12(\DT m_{\hat\lambda}-k)\). The map \(\lambda\mapsto\DT(r)\,m_\lambda-k-2a\lambda\) falls strictly and is
nonnegative at \(\lb\), because \(r\ge r_{\lb}\). If \(\hat\lambda<\lb\), then
\(a\hat\lambda=\Psi\ge\tfrac12(\DT m_{\hat\lambda}-k)>a\hat\lambda\), a contradiction; so \(\hat\lambda=\lb\). Let
\(s_-=\lb\,\sigma_L(-1)\) be the low-class sale rate. Lemma~\ref{lem:a5-bounds}(c) gives \(\sigma_L(1)=0\), so
\(\mu(-1)=\bar b/(2\bar b+s_-)\), and \(U_L(-1)=\rho\DT\bar b/(2\bar b+s_-)-k\) falls strictly in \(s_-\). If
\(\rho\DT m_{\lb}\ge k\), then \(U_L(-1)>0\) at every \(s_-<\lb\), and (E5) gives \(s_-=\lb\). Otherwise \(U_L(-1)\) is
positive at \(s_-=0\), because \(\rho\DT>2k\), and negative at \(s_-=\lb\), so (E5) admits only the \(s_-\in(0,\lb)\)
at which \(U_L(-1)=0\). So \(s_-=\hat s_-(r)\), and as in the first case the tuple is recruiting research at the cap.

Every equilibrium on \((r_N,r_C]\) is therefore floor-entry research or recruiting research at the cap. By parts (b)
and (c), the equilibria are exactly these two on \((r_N,r_Q)\), and recruiting research at the cap alone on
\([r_Q,r_C]\). For \(r\in(r_C,r_2]\), \(r>r_U\) gives \eqref{eq:a5-U}, and \(M_{\lb}<\tau(r)\).
Lemma~\ref{lem:a5-strong} then gives a unique equilibrium, with research at the cap, full orders and
\(\zeta(P(1))=0\), which is floor-entry research at the cap. This proves part (d).
\end{proof}

The research \(\lambda_Q(\DT(r))\) of floor-entry research and the recruitment intensity \(\lambda_E(r)\) both rise
with strength. On \((r_N,r_C]\), the floor margin has the sign of \(\lambda_Q(\DT(r))-\lambda_E(r)\); it is negative at
\(r_N\), where \(\lambda_Q\) starts from zero, and it rises, so \(\lambda_Q\) crosses the recruitment intensity once, at
\(r_Q\). From \(r_Q\) to \(r_C\), a purchase at the research \(\min\{\lambda_Q,\lb\}\) recruits, and only recruiting
research at the cap remains. Above \(r_C\), even research at the cap leaves the purchase posterior below the
threshold, and entry returns to the floor.
